\scanpage{48}
\diagram{\begin{tikzcd}[ampersand replacement=\&, column sep=small]0\arrow[r,no head]\&P'\arrow[r,no head]\&P\arrow[r,no head]\&P''\arrow[r,no head]\&0\&Q\end{tikzcd}}
\margin{\figure{\begin{tikzpicture}\fill (0,0) ellipse (.03 and .015);\end{tikzpicture}}}
\deleted{$\unsure{nP''=0}$}
\[\text{a)}\quad nQ=Q\qquad\text{b)}\quad {}_nP\unsure{\mathbin{-}|\longrightarrow}P''\quad\text{épimorphisme}\]
Toute biextension \inserted{$E$} de $P\times Q$ \deleted{\unsure{définit}} par $G$ définit

a) Une biextension \inserted{$E'$} de $P'\times Q$ par $G$

d'où un accouplement
\diagram{\begin{tikzcd}[ampersand replacement=\&]{}_nP'\times{}_nQ\arrow[r,"E'(n)"]\&G\end{tikzcd}}
\margin{\figure{\begin{tikzpicture}\draw (-.6,1.5)..controls (-.1,1.5) and (0,1.5)..(0,1.2)--(0,-.7)..controls (0,-.9) and (.1,-.9)..(.7,-.9);\end{tikzpicture}}\deleted{\unsure{Si} la suite $0\longrightarrow P'\longrightarrow P\longrightarrow P''\longrightarrow0$ est \unsure{splittée} \deleted{\unsure{ou}} ou elle est
\[\deleted{\operatorname{Hom}(P\otimes Q,G)\longrightarrow\operatorname{Hom}(P'\otimes Q,G)}\]
\uline{\unsure{loc.}} \uline{\unsure{splittée}} et
\[\underline{\operatorname{Hom}}(\unsure{P'|P''}\otimes Q,G)=0\]}}

b) Une biextension de ${}_nP\times Q$ par $G$, i.e. un accouplement
\diagram{\begin{tikzcd}[ampersand replacement=\&]{}_nP\times{}_nQ\arrow[r,"\varphi"]\&G\end{tikzcd}}

\deleted{\unsure{Donc le}} Ces deux $(E',\varphi)$ satisfont la condition que
\[\varphi|({}_nP'\times{}_nQ)=E'(n)\]

\margin{\figure{\begin{tikzpicture}\fill (0,.45) circle (.02);\fill (0,0) circle (.03);\end{tikzpicture}}}
Un hom de biextensions $E_1\longrightarrow E_2$ définit un hom $(E'_1,\varphi_1)\mathbin{-}(\unsure{E_2|E'_2},\varphi_2)$, i.e. \inserted{(déf.)} un hom de biextensions $u:E'_1\longrightarrow E'_2$, tel que $\varphi\sim$

\uline{Je dis que le foncteur} $E\longmapsto(E',\varphi)$ \uline{est une équiv. de catégories.}

\deleted{\deleted{a) \unsure{pl. fid.}}\quad\deleted{b) \unsure{pl. fid.}}
\[
\deleted{\text{a)}\quad\operatorname{Hom}(P\otimes Q,G)\xrightarrow{\sim}\operatorname{Hom}(P'\otimes Q,G)}
\]
\[\deleted{\text{b)}\quad E_1\simeq E_2\quad\text{ssi}\quad(E'_1,\varphi_1)\simeq(E'_2,\varphi_2)}\]
\uline{Pl. fidèle} signifie aussi que $E\unsure{\vdash|\mapsto}E'(P',Q)$ est pl. fid., ou encore, par la suite exacte, que
\[
(*)\quad\operatorname{Hom}(P''\otimes Q,G)=0,\qquad
(**)\quad\operatorname{Hom}(P\otimes Q,G)\longrightarrow\operatorname{Hom}(P'\otimes Q,G)\quad\uuline{\text{surj.}}
\]
\deleted{\unsure{Ce qui est} \illegible{2}}
\deleted{$(***)$ si $\unsure{E\simeq1}$ tel que $(E',\varphi)$ \illegible{1} \unsure{alors} $\unsure{E\simeq1}$}
\inserted{Si $E$ est tel que $E'\simeq1$, $\varphi=0$, alors $E\simeq1$}

La première assertion résulte de a) et b) qui impliquent $P''\otimes Q=0$; la deuxième résulte de c) [dans la démonstration de c), les deux \unsure{mbs|objets} de $(**)$ sont nuls!].

\deleted{\uline{\unsure{Ess. surj.}}} Pour $(***)$, on note que si $E'=0$, alors par la suite exacte $E$ est déterminé par un élément de $\operatorname{Ext}^1(P''\otimes^LQ,G)$, \unsure{or ceci équivaut à} un élément de $\operatorname{Hom}(P''\otimes_1Q,G)$, on trouve que c'est ce qui définit $\varphi$, on gagne.}

\uline{Corollaire 1} Supposons, pour un $\Pi$ de nbs premiers que
\[\left(\varinjlim_{\substack{n\\\text{\unsure{composé}}\ \Pi}}{}_nP\right)\longrightarrow P''\quad\text{surjectif},\qquad\text{et que}\quad p\in\Pi\Longrightarrow pQ=Q.\]
Alors on a une équivalence $E\longmapsto(E',\varphi)$, où $E'$ est une
\margin{47}
